\section{The exact polynomial and its path partition}\label{lc:source}
We import only the finite operator identification from \cite{FinitePaper}.
The matrix needed in the estimates is explicit. On \(1,z,\ldots,z^d\),
let
\begin{equation}\label{lc:banded}
 \mathsf C z^j=2\cos((2j+1)\theta)z^j
       -s_j^2(x+y)z^{j-1}-s_js_{j-1}xy\,z^{j-2},
 \qquad s_j=2\sin(j\theta).
\end{equation}\aliascurrentequation{\label{sa:operator}}
Negative powers vanish because \(s_0=0\). In the finite Fourier theorem
\(\mathsf C=-T\). To specify the imported identification, let
\(\Omega_s\) be the block of the physical pair operator on
\[
 e_u\otimes e_{[s-u]_N},\qquad 0\le u<N,\qquad 0\le s<N.
\]
Here \(e_0,\ldots,e_d\) is its standard one-strand basis, and
\([v]_N\) is the representative in \(\{0,\ldots,N-1\}\).
For \(y\ne0\), specialize \(x=Q^{s+1}y\), set
\(\zeta_u=-q^{-1}xQ^{-u}\), and define
\(\mathsf E_{uj}=\zeta_u^j\), \(G_{uu}=q^{du}\). Then
\begin{equation}\label{lc:residue-adapter}
 \mathsf E\mathsf C\mathsf E^{-1}
                     =G^{-1}\Omega_s^{\mathsf T}G.
\end{equation}
The \(N\) blocks form the \(N^2\)-dimensional pair space. The identity
includes their nilpotent parts, not just eigenvalues. The face
polynomial retains independent \(x,y\); its specialized values recover
the matrix entries by the adapter in \cite[equation (4.5)]{FinitePaper}.
If \(f(\mathsf C)_{ij}\) is the coefficient from \(z^j\)
to \(z^i\), the homogeneous polynomial is
\[
 H_f(x,b,c,y)=\sum_{0\le i\le j\le d}f(\mathsf C)_{ij}b^ic^{d-j}.
\]
Its entries are polynomials in \(x,y\), and its total face degree is \(d\).

Set
\begin{equation}\label{lc:diagonal}
 \sigma=-2-u,\qquad a_j=2+2\cos((2j+1)\theta),\qquad
 \alpha_j=a_j+u,\qquad A=x+y,\quad P=xy.
\end{equation}
The triangular matrix \(\mathsf M=2I+\mathsf C+uI\) has diagonal
\(\alpha_j\), and lowering entries \(-s_j^2A\), \(-s_js_{j-1}P\).
Its diagonal product is \((-1)^N\chi_N(-2-u)\). Consequently
\begin{equation}\label{lc:adjugate}
 \mathcal P_N(x,b,c,y;u):=(-1)^dH_V(x,b,c,y;-2-u)
       =\sum_{i\le j}(\operatorname{adj}\mathsf M)_{ij}b^ic^{d-j}.
\end{equation}
The identity follows first where \(\mathsf M\) is invertible and then
by polynomial equality, including every singular spectral value.

Expand its triangular inverse by the strictly lowering part and
multiply by all diagonal factors. A path \(j\to i\) uses steps of length
one or two. Its weight is the product of \(s_r^2A\) for a step
\(r\to r-1\), \(s_rs_{r-1}P\) for a step \(r\to r-2\), and
\(\alpha_r\) for precisely the rows it does not visit.
The middle monomial \(b^ic^{d-j}\) completes its term.
This proves the path rule, including the case of a zero unused factor.

Partition \(\mathcal P_N=\mathcal P_{\mathrm C}+\mathcal P_{\mathrm L}
+\mathcal P_{\mathrm R}\) by the sectors in the introduction.
Their common sign relative to \(H_V\) does not affect modulus bounds.
Every sector has spectral degree at most \(d\). At \(u=0\), every
path weight is nonnegative. At odd \(N\), a path missing the middle
row has its unused factor \(a_m=0\).

\begin{figure}[ht]
\centering
\begin{tikzpicture}[x=.82cm,y=.68cm,>=Stealth,font=\small]
 \fill[blue!6] (3.65,-.45) rectangle (4.35,3.5);
 \foreach \j in {0,...,8}{
  \node at (\j,3.8){\(\j\)};
  \foreach \h in {0,1.5,3}{\fill[black!35] (\j,\h) circle(1.1pt);}
 }
 \draw[->,very thick] (3,3)--(2,3);
 \draw[->,very thick] (2,3)--(0,3);
 \node[left] at (-.4,3){left};
 \draw[->,very thick,blue!65!black] (6,1.5)--(4,1.5);
 \draw[->,very thick,blue!65!black] (4,1.5)--(2,1.5);
 \draw[->,thick,dashed,red!65!black] (5,1.7) to[bend right=24] (3,1.7);
 \node[left] at (-.4,1.5){central};
 \draw[->,very thick] (8,0)--(7,0);
 \draw[->,very thick] (7,0)--(5,0);
 \node[left] at (-.4,0){right};
 \node[below,align=center] at (4,-.6){middle row \(m=4\)\\
           unused factor \(u\) on the dashed jump};
\end{tikzpicture}
\caption{The partition at the illustrative odd level \(N=9\).
A central path may visit row \(m\), or skip it by \(m+1\to m-1\).
The skipped row contributes its unused factor \(u\), retained in the
exact central polynomial even though it vanishes at \(u=0\).}
\label{lc:path-figure}
\end{figure}


\subsection{Cleared cofactors}
Let \(R_j(b;u)\) be the determinant obtained from the block of
\(\mathsf M\) on rows \(0,\ldots,j\) by replacing its last row with
\((1,b,\ldots,b^j)\). Its path expansion gives
\begin{equation}\label{lc:R-recurrence}
\begin{split}
 R_0&=1,\qquad R_{-1}=R_{-2}=0,\\
 R_j(b;u)&=b^j\prod_{r<j}\alpha_r+s_j^2A R_{j-1}(b;u)
                +s_js_{j-1}\alpha_{j-1}P R_{j-2}(b;u).
\end{split}
\end{equation}\aliascurrentequation{\label{sa:R}}
Equivalently, group the paths by their starting row:
\begin{equation}\label{sa:path}
 \begin{gathered}
 R_j=\sum_{i=0}^j b^i\prod_{r<i}\alpha_r K_{i,j},\qquad
 K_{i,i}=1,\quad K_{i,i-1}=0,\\
 K_{i,k}=s_k^2A K_{i,k-1}+s_ks_{k-1}\alpha_{k-1}P K_{i,k-2}.
 \end{gathered}
\end{equation}
The interval polynomial $K_{i,j}$ sums the paths between rows $i$ and $j$.
Both this formula and \eqref{lc:R-recurrence} are polynomial identities
when a diagonal factor vanishes.

The first term at \(j=0\) is the empty product one. Each \(R_j\)
has face degree \(j\) and spectral degree at most \(j\). Put
\[
 D_b=R_{m-1}(b;u),\qquad E_b=\alpha_{m-1}R_{m-2}(b;u),
\]
with \(E_b=0\) when \(m=1\). Subscripts \(b,c\) here denote substitution
of the middle variable, not differentiation.

\begin{lemma}[Exact central factors]\label{lc:central-factors}
At \(N=2m\),
\begin{equation}\label{lc:even-factor}
 \mathcal P_{\mathrm C}
 =s_m^2A D_bD_c+s_ms_{m-1}P(D_bE_c+E_bD_c).
\end{equation}\aliascurrentequation{\label{sa:even}}
At \(N=2m+1\), with \(G_b=R_m(b;u)\),
\begin{equation}\label{lc:odd-factor}
 \mathcal P_{\mathrm C}=G_bG_c+u s_m^2P D_bD_c.
\end{equation}\aliascurrentequation{\label{sa:odd}}
Both are polynomial identities for every \(u\).
\end{lemma}
\begin{proof}
An even central path crosses the cut exactly once, by \(m\to m-1\),
\(m+1\to m-1\), or \(m\to m-2\). Their weights give the three
terms. Reflection \(r\mapsto d-r\) identifies upper and lower
cofactors: \(a_{d-r}=a_r\), and the reflected step products agree.
At odd level, paths visiting \(m\) split there to give \(G_bG_c\).
A path skipping it has the unique jump \(m+1\to m-1\), with weight
\(s_{m+1}s_mP=s_m^2P\); its unused diagonal is \(u\).
This proves the second term without dividing by \(u\).
\end{proof}

The same reflection interchanges the left and right sectors and \(b,c\).
A one-step at \(r\) reflects to one at \(N-r\), with \(s_{N-r}=s_r\);
a two-step reflects to one at \(N+1-r\), with product \(s_{r-1}s_r\).
It therefore suffices to estimate the left sector.

